\SetPageTheme{challenge}
\PageHeader{Provocare Finala}{Aplicam independent tot traseul}{Final}

\begin{bossbox}[Problema de consolidare]
  Scrie un program care citeste un numar natural si afiseaza:
  \begin{itemize}
    \item cate cifre are;
    \item suma cifrelor sale;
    \item cea mai mare cifra;
    \item daca este sau nu palindrom.
  \end{itemize}
  Explica pe scurt ce rol are fiecare variabila folosita.
\end{bossbox}

\begin{guidebox}[Criteriu de reusita]
  Daca rezolvi aceasta problema fara sa schimbi schema principala a buclei, inseamna ca ai inteles cu adevarat mecanismul editiei.
\end{guidebox}

\newpage
\SetPageTheme{paper}
\pagecolor{white}
\color{black}
\thispagestyle{empty}
\null
\clearpage

\SetPageTheme{challenge}
\pagecolor{cqbg}
\color{cqtext}
\thispagestyle{empty}
\begin{center}
\vspace*{2.8cm}
{\Huge\bfseries \IssueTitle}\par
\vspace{0.6em}
{\Large \IssueSubtitle}\par
\vspace{2em}

\begin{tcolorbox}[colback=cqcodebg,colframe=cqcyan,boxrule=1pt,arc=2mm,width=0.82\linewidth]
\small
\textbf{Ce ramane dupa aceasta editie}
\begin{itemize}
\item calculatorul nu citeste dintr-o privire;
\item el descompune si verifica;
\item ordinea informatiei conteaza;
\item algoritmii siguri lucreaza in pasi mici si controlabili.
\end{itemize}
\end{tcolorbox}

\vfill
{\ttfamily\small Code Quest \textbullet{} Volume 2 \textbullet{} Issue 1}
\vspace*{2.4cm}
\end{center}
